\chapter{Alur End-to-End}
\label{ch:08-end-to-end}

\section{Satu Perintah dari Klik sampai Baris Basis Data}

\figref{fig:e2e-sequence} menyatukan seluruh bab sebelumnya menjadi satu urutan
pesan. Yang digambarkan adalah kasus yang paling banyak melibatkan komponen, yaitu
pemberian pakan manual yang berhasil, karena kasus itu menyentuh kelima bagian
sistem: peramban, aplikasi dan basis data, broker, serta perangkat.

\figureasset[\linewidth]{e2e-sequence}{Urutan ujung-ke-ujung satu perintah pakan manual: dari
klik pengguna, penulisan baris basis data, publikasi perintah, eksekusi gate,
sampai pembaruan status perintah dan riwayat pakan.}{fig:e2e-sequence}

Berikut yang terjadi pada sumbu waktu, dimulai dari klik. Pengguna menekan tombol
pemberian pakan dan mengonfirmasi jumlah porsi pada dialog; aplikasi menyusun
pengenal perintah baru berbentuk \code{cmd-} diikuti uuid, lalu membuka transaksi
tulis dan menyisipkan dua baris sekaligus, yaitu satu baris pada tabel
\code{commands} berstatus terkirim dan satu baris pada tabel \code{feeding_logs}
berstatus menunggu dengan level sebelum diambil dari telemetri terakhir. Transaksi
ditutup, dan baru setelah itu pesan JSON dipublikasikan ke topik perangkat dengan
QoS~1. Broker meneruskan pesan ke perangkat; perangkat memeriksa topik, memasukkan
payload ke antrean, dan pada putaran \code{loop()} berikutnya mengambil satu
perintah dari antrean. Karena gate masih idle, servo siap, dan level pakan tidak
nol, perangkat memulai porsi pertama: gate terbuka \PortionOpenMs{} ms, ditutup,
diberi jeda \SettleMs{} ms, lalu jeda \GapMs{} ms sebelum porsi berikutnya.
Setelah porsi terakhir, perangkat mengukur ulang jarak, mempublikasikan telemetri,
lalu mempublikasikan ack berisi pengenal perintah yang sama dan hasil \code{ok}.
Ack diterima worker, yang memanggil handler ack; handler menandai baris perintah
sukses beserta waktu ack, menandai baris riwayat pakan sukses, dan mengisi level
sesudah dari baris telemetri terakhir. Pada saat itulah baris basis data yang
dibuat di awal berubah menjadi hasil akhir, dan halaman Status, History, serta
Control menampilkannya pada penyegaran berikutnya.

Bila salah satu mata rantai itu putus, urutannya berhenti pada titik yang
berbeda dan hasilnya tetap terlihat. Broker yang tidak meneruskan pesan membuat
perintah tetap berstatus terkirim sampai sweeper menandainya \code{timeout}.
Perangkat yang menolak karena keadaan tertentu menjawab dengan hasil selain
\code{ok}, sehingga perintah ditandai gagal dan notifikasi pemberian pakan gagal
muncul. Perangkat yang mati total mengubah topik status menjadi offline, sehingga
halaman Status menampilkan perangkat tidak terhubung. Tidak ada jalur yang
berakhir tanpa status.

\section{Matriks Pesan MQTT dan Aksi Basis Data}

\tabref{tab:08-matriks} memetakan setiap jenis pesan yang mungkin masuk ke aksi
basis data yang ditimbulkannya. Tabel ini adalah ringkasan dari empat handler pada
\file{handlers.py} dan satu thread sweeper pada \file{runtime.py}.

\begin{table}[H]
\centering
\caption{Matriks pesan masuk terhadap aksi basis data.}
\label{tab:08-matriks}
\small
\begin{tabular}{@{}p{3.0cm}p{5.2cm}p{3.0cm}p{3.2cm}@{}}
\toprule
\textbf{Pesan} & \textbf{Aksi basis data} & \textbf{Tabel} & \textbf{Pemicu
lanjutan} \\
\midrule
\code{telemetry} & Sisip satu baris telemetri, perbarui waktu terakhir terlihat,
tutup notifikasi offline, evaluasi ambang pakan. & \code{telemetry}, \code{devices},
\code{notifications} & Notifikasi pakan menipis atau habis. \\
\code{status} online & Tandai perangkat online, tutup notifikasi offline. & \code{devices},
\code{notifications} & -- \\
\code{status} offline & Tandai perangkat offline, buat notifikasi offline bila
belum ada. & \code{devices}, \code{notifications} & Notifikasi perangkat offline. \\
\code{ack} dengan hasil ok & Tandai perintah sukses, tandai riwayat pakan sukses,
isi level sesudah. & \code{commands}, \code{feeding_logs} & -- \\
\code{ack} dengan hasil lain & Tandai perintah dan riwayat pakan gagal, buat
notifikasi kegagalan. & \code{commands}, \code{feeding_logs}, \code{notifications} & Notifikasi
pemberian pakan gagal. \\
\code{state} event jadwal & Sisip satu baris riwayat terjadwal tanpa pengenal
perintah, isi level sebelum dan sesudah. & \code{feeding_logs}, \code{devices} & Tampil di
halaman History sebagai terjadwal. \\
Tanpa ack dalam \CmdTimeoutS{} detik & Tandai perintah timeout, tandai riwayat
pakan gagal dengan pesan timeout. & \code{commands}, \code{feeding_logs} & -- \\
Tanpa data dalam \OfflineS{} detik & Tandai perangkat offline, buat notifikasi
offline. & \code{devices}, \code{notifications} & Notifikasi perangkat offline. \\
\bottomrule
\end{tabular}
\end{table}

\section{Ketertelusuran Kebutuhan}

\tabref{tab:08-trace} menautkan setiap kebutuhan fungsional dan non-fungsional ke
bab yang membahasnya, gambar yang menggambarkannya, dan berkas yang
mengimplementasikannya. Tabel ini dipakai sebagai daftar periksa: tidak ada
kebutuhan yang boleh hanya muncul di dokumen tanpa jejak implementasi, dan tidak
ada perilaku penting yang hanya ada di kode tanpa penjelasan.

{\scriptsize
\begin{longtable}{@{}p{1.1cm}p{2.9cm}p{3.1cm}p{7.3cm}@{}}
\caption{Ketertelusuran kebutuhan ke bab, gambar, dan berkas implementasi.}
\label{tab:08-trace} \\
\toprule
\textbf{Kode} & \textbf{Kebutuhan} & \textbf{Bab dan gambar} & \textbf{Berkas
implementasi} \\
\midrule
\endfirsthead
\toprule
\textbf{Kode} & \textbf{Kebutuhan} & \textbf{Bab dan gambar} & \textbf{Berkas
implementasi} \\
\midrule
\endhead
\bottomrule
\endfoot
FR-01 & Grafik sisa pakan periodik untuk rentang 1, 6, atau 24 jam &
\chapref{ch:05-website}, \figref{fig:web-flowchart} & \file{live.py},
\file{data.py}, \file{db.py} (\code{telemetry_series}) \\
FR-02 & Tampilan keadaan terakhir perangkat dengan penyegaran otomatis dan penanda
kondisi pakan & \chapref{ch:05-website}, \tabref{tab:05-halaman} &
\file{status.py}, \file{schema.sql}
(\code{v_device_status}) \\
FR-03 & Pemberian pakan manual dari dashboard & \chapref{ch:05-website},
\figref{fig:web-seq-manual-feed-ok}, \figref{fig:web-seq-manual-feed-fail} &
\file{commands.py}, \file{handlers.py},
\file{esp32_feeder.ino} \\
FR-04 & Permintaan keadaan perangkat saat ini & \chapref{ch:06-protokol-mqtt},
\tabref{tab:06-command} & \file{control.py},
\file{commands.py}, \file{protocol.cpp} \\
FR-05 & Pengelolaan jadwal dan pengirimannya ke perangkat &
\chapref{ch:05-website}, \figref{fig:web-seq-schedule-sync} &
\file{schedules.py}, \file{db.py}
(\code{insert_schedule}), \file{esp32_feeder.ino}
(\code{scheduleTick}) \\
FR-06 & Pengaturan ambang notifikasi pakan menipis & \chapref{ch:05-website},
\tabref{tab:06-command} & \file{control.py},
\file{db.py} (\code{set_low_feed_threshold}),
\file{handlers.py} \\
FR-07 & Notifikasi pakan menipis, habis, gagal, dan perangkat offline beserta
penandaan dibaca & \chapref{ch:05-website}, \figref{fig:web-seq-telemetry-alert},
\tabref{tab:07-kamus} & \file{handlers.py},
\file{alerts.py}, \file{schema.sql} \\
FR-08 & Riwayat pemberian pakan per periode dengan level sebelum dan sesudah &
\chapref{ch:05-website}, \figref{fig:fw-seq-scheduled-feed} &
\file{history.py}, \file{db.py}
(\code{list_feeding_logs}) \\
NFR-01 & Rantai perangkat ke MQTT ke basis data ke web berjalan tanpa kehilangan
pesan, dengan penyambungan ulang otomatis & \chapref{ch:02-arsitektur},
\figref{fig:e2e-topology} & \file{mqtt.py}, \file{handlers.py},
\file{net.cpp} \\
NFR-02 & Latensi perintah pada jaringan lokal dan penandaan timeout &
\chapref{ch:05-website}, \tabref{tab:05-sweeper} & \file{runtime.py},
\file{config.py} (\code{COMMAND_TIMEOUT_SEC}) \\
NFR-03 & Akurasi konversi jarak ke persentase sisa pakan & \chapref{ch:03-perangkat-keras}
(bagian kalibrasi hopper) & \file{handlers.py} (\code{compute_pct}),
\file{esp32_feeder.ino} (\code{pctFromCm}) \\
NFR-04 & Akurasi eksekusi jadwal terhadap jam terjadwal & \chapref{ch:04-firmware},
\figref{fig:fw-seq-scheduled-feed} & \file{esp32_feeder.ino}
(\code{scheduleTick}, jendela \WindowS{} detik) \\
NFR-05 & Batasan keamanan: broker tanpa autentikasi dan tanpa TLS, hanya untuk
jaringan lokal & \chapref{ch:02-arsitektur}, \chapref{ch:09-catatan-asumsi} &
\file{broker.py} (autentikasi anonim) \\
NFR-06 & Batasan skala: sistem divalidasi untuk satu perangkat &
\chapref{ch:01-pendahuluan}, \tabref{tab:07-kamus} & \file{schema.sql}
(seed \code{feeder01}), \file{config.h} \\
\end{longtable}
}

\section{Anggaran Waktu Satu Perintah}

Dengan nilai yang berlaku sekarang, satu porsi memerlukan \PortionCycleMs{} ms
ditambah jeda antar porsi sebesar \GapMs{} ms bila porsi lebih dari satu.
Karena itu perintah dua porsi memerlukan $2 \times \PortionCycleMs + \GapMs$
milidetik, dan perintah sepuluh porsi memerlukan
$10 \times \PortionCycleMs + 9 \times \GapMs$ milidetik. Angka-angka itu adalah
batas bawah, bukan batas atas: setelah gate terakhir menutup, perangkat masih
mengukur jarak, mempublikasikan telemetri, lalu mengirim ack. Batas waktu
\CmdTimeoutS{} detik pada platform dipilih jauh di atas total tersebut, sehingga
perintah sepuluh porsi pun tidak akan ditandai timeout ketika perangkat sehat.

Untuk perintah yang tidak menggerakkan gate, seperti pengaturan ambang atau
permintaan keadaan, waktu proses di perangkat hanya satu putaran \code{loop()},
dan round-trip ditentukan oleh latensi jaringan lokal. Angka itu tidak dijanjikan
di dokumen ini karena tidak diukur secara sistematis; yang dijamin oleh rancangan
adalah batas atas \CmdTimeoutS{} detik untuk setiap perintah, apa pun tipenya.
